% FILE: 01_research_question.tex
% Project: research-open-ontologies
% Role: open-ontology-research-and-rdf-authority

\section{Research Question}
\label{sec:research-open-ontologies-question}

\subsection{Problem Statement}
\label{sec:research-open-ontologies-problem}

Process intelligence claims --- fitness scores, conformance verdicts, M\&A board projections,
lifecycle governance decisions --- are only as trustworthy as the ontology layer that defines
what each term means and what evidence is required to assert it. Prior work in process mining
produces claims in natural language or ad hoc data structures without a machine-readable,
dependency-validated knowledge graph that an auditor can inspect, replay, or challenge.

The specific problem this project addresses is: \emph{how can a distributed RDF corpus, spanning
multiple research sub-programs and using both public standard vocabularies and local extension
namespaces, be validated as acyclic, dependency-consistent, and queryable without requiring all
TTL files to live in a single directory?} A secondary problem is how that corpus can then
serve as the input to a code manufacturing pipeline that emits type-safe, receipted artifacts
rather than ad hoc text.

\subsection{Old-Regime Assumption Challenged}
\label{sec:research-open-ontologies-old-regime}

The old regime assumes that process intelligence claims are self-evident from algorithm output.
A fitness score of 0.97 from a token replay is taken as proof of conformance without asking: What
named witness grounds this threshold? What ontology defines \texttt{fitness} as a typed property?
What receipt chain connects this score to the event log that produced it? What LTL safety
invariant prevents a non-compliant process from being actuated?

The assumption challenged here is that natural-language documentation of process evidence is
sufficient for board-level or regulatory purposes. This project asserts the contrary: every
claim must be backed by a machine-readable ontology triple, a named witness marker, and a
BLAKE3-receipted artifact. A claim that cannot be expressed as a SPARQL query against a
validated RDF graph is not an admissible process intelligence claim under this program's
COVENANT.

\subsection{Hypothesis}
\label{sec:research-open-ontologies-hypothesis}

The hypothesis is that a federated RDF corpus --- distributed across multiple sub-programs,
validated for dependency acyclicity and triple consistency, and grounded in public standard
vocabularies --- can serve as a reliable authority surface for manufacturing board-admissible
process intelligence artifacts via SPARQL-to-Tera pipelines.

Specifically: if the federated dependency graph reports \texttt{ACYCLIC}/\texttt{CONSISTENT}
with zero circular dependencies and zero dangling references, then SPARQL queries against the
merged graph will return results that can be transformed via Tera templates into typed Rust
enums, WIT boundary definitions, and Markdown documentation --- each with a verifiable BLAKE3
receipt. The hypothesis is PARTIAL at the time of this writing because the manufacturing
pipeline has not yet been validated end-to-end against a receipted ALIVE checkpoint for the
ontology project itself (as distinct from the parent program checkpoint).

\subsection{Success Criteria}
\label{sec:research-open-ontologies-success}

Success for this project is defined by the following gates, derived from the parent program
COVENANT and the Van der Aalst Constitution (process-mining-chicago-tdd):

\begin{enumerate}
  \item The federated dependency graph reports \texttt{gate\_pass: true} with zero circular
        dependencies and zero dangling references. \textbf{[MET --- Phase~4, 2026-06-01]}
  \item All TTL files in the corpus load successfully via rdflib with no parse errors.
        \textbf{[MET --- 14/14 files, Phase~7 roundtrip report]}
  \item SPARQL smoke queries return non-empty results against the unified graph.
        \textbf{[PARTIALLY MET --- 4/4 queries pass, but \texttt{count\_checkpoints=0} and
        \texttt{count\_artifacts=0} due to namespace prefix mismatch]}
  \item TTL syntax validation via \texttt{rapper}/\texttt{librdf-utils} passes for all 23
        TTL files. \textbf{[NOT MET --- tooling not installed; Gate~3 FAIL in conformance audit]}
  \item A receipt TTL exists in the ontology corpus providing RDF-native audit trail for the
        ontology layer itself. \textbf{[NOT MET --- RECEIPT\_NOT\_FOUND in Phase~7 report]}
  \item The ggen manufacturing pipeline emits at least one receipted artifact from a SPARQL
        query against this corpus. \textbf{[PARTIALLY MET --- ggen.toml rules defined,
        receipts present in \texttt{.ggen/receipts/}, ALIVE reissue blocked pending
        Gates~3 and~4]}
\end{enumerate}
